\documentclass[]{article}
\usepackage{color}
\usepackage{amsfonts}
%opening
\title{Scientific project}
\author{}

\begin{document}

\maketitle

\section{Introduction}

In computer science, a \textit{decision procedure} is an algorithm that, given a logical formula $\varphi$, determines whether $\varphi$ can be made true (that is, is \textit{satisfiable}). Depending on the formula, such an algorithm may exist or not. A class of formulas for which a decision procedure exists is a \textit{decidable logical fragment}.

Decision procedures are used as black boxes in many areas of computer science. In software verification, a common approach is to encode the behaviour of a computer system as a logical formula. Verifying critical properties of the system then amounts to studying the satisfiability of the formula by means of a decision procedure. Other applications include powerful software called \textit{SMT solvers}, which are able to study the satisfiability of highly expressive formulas from various fragments thanks to a combination of distinct reasoning tools, including decision procedures.

Many decision procedures rely on the concept of \textit{finite-state automata}. A finite-state automaton is a mathematical structure that is syntactically simple, easy to manipulate algorithmically yet semantically highly expressive. For this reason, automata often serve as practical data structures in decision procedures. The efficiency of these so-called \textit{automata-based decision procedures} strongly depends on the type of automata that they use, and also on syntactical restrictions that can be made on them. Informally, our research project aims at identifying restricted classes of automata that could be manipulated more efficiently in the context of automata-based decision procedures.

\section{State of the art}

Automata-based decision procedures work as follows. Given a formula $\varphi$, the idea is to build an automaton representing $\varphi$ (that is, accepting exactly the encoding of its models) by first associating each atomic sub-formula of $\varphi$ to an automaton, and then by following the syntactical structure of $\varphi$ and combining those automata accordingly. The last step is to check that the automaton obtained in this way accepts a non-empty language. 

This approach was first used by Büchi to decide the first-order additive arithmetic of integer numbers (also known as Presburger arithmetic)\cite{buc60} and the monadic second-order theory of one successor~\cite{buc62}. The latter logical fragment requires to use infinite-word automata, for instance Büchi automata. The seminal contributions of Büchi were then extended in various directions. For instance, a decision procedure for the first-order additive arithmetic with integer and real variables (an extension of Presburger arithmetic) by means of Büchi automata was developed in~\cite{}. The monadic second-order theory of order (that can be regarded as an extension of S1S on arbitrary domains) was proved equivalent to an advanced form of automata (called automata on linear orderings) on scattered countable domains~\cite{}.

These automata-based decision procedures are interesting in themselves, since they establish strong links between formal language theory and logic. In addition, they provide several advantages compared to other types of decision procedures (for instance, quantifier elimination). Automata provide a symbolic representation of all models of a formula. This makes it possible to search for a model satisfying some optimization criteria, as well as to compute the set of models once and then query it multiple times for different properties of interest. Moreover, even in the cases where the worst-case complexity of these methods seems prohibitive (deciding Presburger arithmetic with automata is triply-exponential in time), their practical implementations (for instance LASH~\cite{} and Walnut~\cite{} for arithmetic) run successfully on many practical instances.

The main challenge in representing a formula by an automaton is to translate the effect of universal quantifiers. Classically, one circumvents this difficulty by means of the identity $\forall=\neg\exists\neg$. The effect of a universal quantifier can thus be translated into a complementation of automaton, followed by a projection, followed itself by a second complementation. Unfortunately, complementing a (non-deterministic) finite-word automaton or a Büchi automaton is a difficult operation. In general, it leads to an exponential blow-up in the number of states of the automaton. 

However, it has been shown recently in~\cite{} that for Büchi automata encoding models of the monadic first-order theory of order on $\mathbb{N}$, it is possible to apply universal quantification directly, without complementation algorithm. The resulting blow-up in the number of states is only exponential, while the classical method incurs at least two exponential blow-ups in general (one for each complementation). A similar observation has been made in the context of quantified integer programming~\cite{}.

\section{Research question}

The research project consists in identifying a syntactical subclass of finite-word automata for which a universal quantification algorithm leads to a single exponential blow-up in the number of states of the automaton. The same question can be asked for Büchi automata. If such a subclass can be characterized, another research objective is to determine which logical fragment can be decided thanks to this subclass of automata. Answering this question would immediately yield a more efficient decision procedure for this logical fragment. Moreover, it would probably highlight some new deep connections between logic and formal language theory. On a practical level, identifying syntactical restrictions under which an operation can be computed more efficiently allows to design new data structures that exploit these restrictions.

Results of this kind have already been obtained in the past. For instance, as shown in~\cite{}, any formula in the first-order additive arithmetic with integer and real variables can be represented by a weak Büchi automaton, that is, by a restricted form of Büchi automaton that can be manipulated more efficiently. This result is interesting in itself since it refines the link between arithmetical theories and automata. In addition, it allows to decide linear arithmetic more efficiently, through the use of data structures tailored to weak Büchi automata. 

One possible way of addressing our research question is probably to opt for an algebraic approach. Finite-state automata have deep connections with finite monoids. For instance, both formalism are equivalent in the case of finite words~\cite{}. One way of showing that finite-word automata can be complemented is to exploit this equivalence, together with the fact that complementing a monoid is a trivial operation. As hinted above, universal quantification is strongly linked to complementation. It is thus likely that regarding the problem as an algebraic one yields valuable insights.

\bibliography{refs}
\end{document}
